\documentclass[10pt,a4paper]{amsart}
\usepackage[margin=19mm]{geometry}
\usepackage{amsmath,amssymb,mathtools}
\usepackage[scr=rsfs,cal=euler]{mathalfa}
\usepackage{xfrac}
\usepackage[unicode,hidelinks,bookmarks=false]{hyperref}
\newcommand{\ie}{i.e.\ }
\newcommand{\cf}{cf.\ }
\newcommand{\resp}{respectively\ }
\newcommand{\Subsection}{\subsection}
\providecommand{\nobreakdash}{\nobreak\hbox{-}}
\setlength{\emergencystretch}{2em}
\multlinegap0pt
\usepackage{bm}
\usepackage{bbm}
\usepackage{dsfont}
\usepackage{xspace}
\usepackage{cancel}
\usepackage{mathtools}
\usepackage{amsthm}
\makeatletter
\@ifundefined{theorem}{\newtheorem{theorem}{Theorem}}{}
\@ifundefined{proposition}{\newtheorem{proposition}[theorem]{Proposition}}{}
\makeatother
\DeclareMathOperator{\Cay}{Cay}
\title[Abelian groups without 3-chromatic Cayley graphs]{Abelian groups without 3-chromatic Cayley graphs}
\author{Mike Krebs, Maya Sankar}
\date{Source: arXiv:2410.11028v2 (CC BY 4.0)}
\begin{document}
\maketitle

\noindent\emph{Source and licence.} Abelian groups without 3-chromatic Cayley graphs. Version arXiv:2410.11028v2 (CC BY 4.0), distributed under the Creative Commons Attribution 4.0 International licence, as recorded in the arXiv licence metadata for this exact version and shown on the abstract page https://arxiv.org/abs/2410.11028v2. Full rights notice: licenses/arxiv-2410.11028-CC-BY-4.0.txt.\par\smallskip

\noindent\textit{Editorial scope.} Counted unit: the Proposition whose source label is \texttt{prop:pi1-for-abelian-Cayley-graph} in the article (source0.tex lines 277--279, source label \texttt{prop:pi1-for-abelian-Cayley-graph}), delivered together with the article's own complete original proof of it (source0.tex lines 281--300, one \texttt{proof} environment of 20 lines). No mathematics is written, repaired or re-derived: statement and proof are the article's own text line for line. Required context retained uncounted: none. Every definition, display and label of those lines is retained unchanged. Omitted from this package and not redistributed: 1--276, 280--280, 301--651. Nothing inside a retained excerpt is omitted or altered: every retained body is delivered exactly as the article's lines are written, so no omission receipt of any kind accompanies this package. Citations: the retained excerpts carry no citation command at all, so no bibliography entry of the article is transcribed and no reference is redistributed. Reference adaptation: none was needed and none was made; every reference inside the delivered unit resolves inside it, so no unresolved reference or citation is produced. Printed slips kept literal: no printed word, symbol, hypothesis or number of the counted statement or of its proof was altered. Layout and class: the article's own class is \texttt{amsart} with packages including amsmath, amssymb, amsthm, fullpage, verbatim, cite, hyperref, cleveref, enumitem, xcolor, tikz, caption, tikz-cd; the wrapper is the lane's pinned \texttt{amsart} editorial wrapper at 10pt on A4, which restates the theorem-like environments the retained text needs and adds the article's own macro definitions that the retained text needs (\texttt{\textbackslash Cay}), with extra wrapper packages (bm, bbm, dsfont, xspace, cancel, mathtools). The delivered numbering is the unit's own. Assets: no retained span contains a figure environment, a graphics-inclusion command, a PDF-page inclusion, a table or a TikZ picture; no image file is copied into this package, the package declares no asset dependency and no image file is redistributed with it. Licence: version arXiv:2410.11028v2 of this article is distributed under the Creative Commons Attribution 4.0 International licence, as recorded in the arXiv licence metadata for this exact version and shown on the abstract page https://arxiv.org/abs/2410.11028v2. A full rights notice is delivered at licenses/arxiv-2410.11028-CC-BY-4.0.txt. Characters: every retained excerpt is pure ASCII, so the delivered engine, XeLaTeX with its default Latin Modern text font, reports no missing character.
\begin{proposition}\label{prop:pi1-for-abelian-Cayley-graph}
Suppose $G$ is an abelian group and $S\subset G$ a symmetric subset that generates $G$. Then $\pi_1(\Cay(G,S))\cong R[S]/H[S]$.
\end{proposition}

\begin{proof}
In this proof, it is convenient to denote a closed walk by the generators used to take each step rather than by the vertices visited. Given $s_1,\ldots,s_\ell\in S$, let $(v_0;s_1,\ldots, s_\ell)$ denote the walk $v_0\dots v_\ell$ of length $\ell$ with $v_i:=v_0+s_1+\cdots+s_i$, and let $[v_0;s_1,\ldots,s_\ell]$ denote the homotopy equivalence class of this walk. 
Observe that $(v_0;s_1,\dots, s_\ell)$ is homotopy equivalent to $(v_0;s_1,\dots,s_{i-1},s_{i+1},s_i,s_{i+2},\dots,s_\ell)$ via a single substitution. That is, permuting the generators $s_1,\ldots,s_\ell$ has no effect on the homotopy equivalence class of the walk $(v_0;s_1,\ldots,s_\ell)$.
Additionally, if $s_{i+1}=-s_i$ then $(v_0;s_1,\dots,s_\ell)$ is homotopy equivalent to $(v_0;s_1,\ldots,s_{i-1},s_{i+2},\dots,s_\ell)$. Because we may freely reorder the generators, it follows that homotopy equivalence is preserved by deleting any two generators satisfying $s_i=-s_j$.

Let $\phi:R[S]\to\pi_1(\Cay(G,S),0)$ be the map that takes $a_1\cdot s_1+\cdots+a_k\cdot s_k\in R[S]$ to $[0;\pm s_1,\ldots,\pm s_1,\ldots,\pm s_k,\ldots,\pm s_k]$, which features either $s_i$ repeated $a_i$ times (if $a_i>0$) or $-s_i$ repeated $-a_i$ times (if $a_i<0$). We claim that $\phi$ is a group homomorphism. If $R_1=a_1\cdot s_1+\cdots+a_k\cdot s_k$ and $R_2=b_1\cdot s_1+\cdots +b_k\cdot s_k$ are elements of $R[S]$ then $\phi(R_1)*\phi(R_2)$ takes the form
$$\phi(R_1)*\phi(R_2)=[0;\pm s_1,\ldots,\pm s_1,\pm s_k\ldots,\pm s_k,\pm s_1,\ldots,\pm s_1,\pm s_k\ldots,\pm s_k].$$
By reordering the generators and canceling pairs of the form $(s_i,-s_i)$, we may write
\[
\phi(R_1)*\phi(R_2)=[0;\pm s_1,\ldots,\pm s_1,\ldots,\pm s_k,\ldots,\pm s_k],
\]
featuring $|a_i+b_i|$ copies of either $s_i$ (if $a_i+b_i>0$) or $-s_i$ (if $a_i+b_i<0$). This is exactly $\phi(R_1+R_2)$.

Additionally, $\phi$ is surjective, as $[0;s_1,\ldots,s_\ell]$ is the image of $1\cdot s_1+\cdots+1\cdot s_\ell$. Thus, $\pi_1(\Cay(G,S),0)\cong R[S]/\ker\phi$. It holds that
\[
\ker\phi=\{(1\cdot s_1+\cdots+1\cdot s_\ell)-(1\cdot s'_1+\cdots+1\cdot s'_m) \mid [0;s_1,\ldots,s_\ell]=[0;s'_1,\ldots,s'_m]\}
\]
is generated by differences of the form $(1\cdot s_1+\cdots +1\cdot s_\ell)-(1\cdot s'_1+\cdots+1\cdot s'_m)$ where $W=(0;s_1,\ldots,s_\ell)$ and $W'=(0;s'_1,\ldots,s'_m)$ differ by a single insertion, deletion, or substitution. If $W$ and $W'$ differ by a single insertion or deletion, the resulting difference is of the form $\pm(1\cdot s+1\cdot (-s))$. If $W$ and $W'$ differ by a single substitution, the resulting difference is of the form $1\cdot s_i+1\cdot s_{i+1}-1\cdot s'_i-1\cdot s'_{i+1}$, where these four generators satisfy $s_i+s_{i+1}=s'_i+s'_{i+1}$.
These are exactly the generators of $H[S]$, so $\ker\phi=H[S]$ and $\pi_1(\Cay(G,S),0)\cong R[S]/\ker\phi\cong R[S]/H[S]$.
\end{proof}

\end{document}
